\section{Задание}

Разработать программное обеспечение для визуализации функций распределения и плотности распределения для следующих типов вероятностных распределений:
\begin{itemize}
    \item Равномерное распределение
    \item Пуассоновское распределение
    \item Экспоненциальное распределение
    \item Нормальное распределение
    \item Распределение Эрланга
\end{itemize}

Программа должна позволять интерактивно изменять параметры распределений и визуализировать графики в реальном времени.

\section{Теоретические сведения}

\subsection{Равномерное распределение}

Равномерное (прямоугольное) непрерывное распределение описывается плотностью вероятности:

\begin{equation}
f(x) = \frac{1}{b-a}, \quad x \in [a,b]
\end{equation}

где $a$ --- нижняя граница интервала, $b$ --- верхняя граница интервала.

Функция распределения имеет вид:

\begin{equation}
F(x) = \begin{cases}
0, & x \leq a \\
\frac{x-a}{b-a}, & a < x \leq b \\
1, & x > b
\end{cases}
\end{equation}

Применяется для моделирования событий с равной вероятностью на заданном интервале, генерации псевдослучайных чисел.

\subsection{Пуассоновское распределение}

Пуассоновское дискретное распределение описывается функцией вероятности:

\begin{equation}
P(X=k) = \frac{\lambda^k e^{-\lambda}}{k!}, \quad k \in \mathbb{N}_0
\end{equation}

где $\lambda > 0$ --- интенсивность (среднее число событий), $k$ --- число событий.

Функция распределения:

\begin{equation}
F(k) = \sum_{i=0}^{k} \frac{\lambda^i e^{-\lambda}}{i!}
\end{equation}

Применяется для моделирования числа событий, происходящих за фиксированный промежуток времени при известной средней частоте (например, число заявок в системе обслуживания).

\subsection{Экспоненциальное распределение}

Экспоненциальное непрерывное распределение характеризуется плотностью:

\begin{equation}
f(x) = \lambda e^{-\lambda x}, \quad x \geq 0
\end{equation}

где $\lambda > 0$ --- интенсивность (скорость наступления события).

Функция распределения:

\begin{equation}
F(x) = 1 - e^{-\lambda x}, \quad x \geq 0
\end{equation}

Применяется для моделирования времени до наступления события в пуассоновском потоке, времени безотказной работы, интервалов между поступлениями заявок.

\subsection{Нормальное распределение}

Нормальное (гауссовское) непрерывное распределение описывается плотностью:

\begin{equation}
f(x) = \frac{1}{\sigma\sqrt{2\pi}} e^{-\frac{(x-\mu)^2}{2\sigma^2}}
\end{equation}

где $\mu$ --- математическое ожидание (среднее значение), $\sigma > 0$ --- стандартное отклонение.

Функция распределения выражается через функцию ошибок и не имеет аналитического выражения в элементарных функциях.

Применяется для моделирования случайных величин, подверженных влиянию множества независимых факторов (по центральной предельной теореме), а также в статистическом анализе.

\subsection{Распределение Эрланга}

Распределение Эрланга является частным случаем гамма-распределения с целым параметром формы. Плотность вероятности:

\begin{equation}
f(x) = \frac{\lambda^k x^{k-1} e^{-\lambda x}}{(k-1)!}, \quad x \geq 0
\end{equation}

где $k \in \mathbb{N}$ --- параметр формы (число фаз), $\lambda > 0$ --- интенсивность (скорость).

Функция распределения выражается через неполную гамма-функцию.

Применяется для моделирования времени обслуживания, состоящего из $k$ последовательных экспоненциально распределенных фаз. Широко используется в теории массового обслуживания.

\section{Описание реализации}

\subsection{Используемые библиотеки}

Программа реализована на языке Python с использованием следующих библиотек:

\begin{itemize}
    \item \textbf{PyQt6} --- создание графического интерфейса пользователя
    \item \textbf{NumPy} --- численные вычисления и работа с массивами
    \item \textbf{SciPy} --- вычисление функций распределения и плотности вероятности
    \item \textbf{Matplotlib} --- построение графиков и рендеринг LaTeX-формул
\end{itemize}

\newpage
\subsection{Метод решения}

Для каждого распределения реализованы следующие этапы:

\begin{enumerate}
    \item Определение диапазона значений аргумента на основе параметров распределения
    \item Генерация массива точек (1500 точек для непрерывных распределений)
    \item Вычисление значений функции плотности и функции распределения с использованием библиотеки SciPy
    \item Построение графиков
\end{enumerate}

Для распределения Пуассона используется дискретный набор значений от 0 до $\lambda + 4\sqrt{\lambda}$.

\section{Вывод}

В ходе выполнения лабораторной работы было разработано программное обеспечение для визуализации и анализа пяти основных типов вероятностных распределений.

Программа обеспечивает интерактивное взаимодействие с параметрами распределений, автоматическую валидацию входных данных и визуализацию графиков функций распределения и плотности в едином окне. Реализованный интерфейс позволяет наглядно изучать влияние параметров на форму распределений, что может быть использовано в образовательных целях и при проектировании систем моделирования.